\documentclass[11pt,a4paper]{article}
\usepackage[utf8]{inputenc}
\usepackage[T1]{fontenc}
\usepackage[catalan]{babel}
\usepackage{lmodern}
\usepackage[sfdefault,lining]{FiraSans}
\usepackage[a4paper,top=2.3cm,bottom=2.2cm,left=1.8cm,right=1.8cm,headheight=1.3cm,headsep=0.4cm,footskip=1cm]{geometry}
\usepackage{amsmath,amssymb}
\usepackage{icomma}
\usepackage{graphicx}
\usepackage[table]{xcolor}
\usepackage{tikz}
\usepackage{fontawesome5}
\usepackage{ccicons}
\usepackage{fancyhdr}
\usepackage{lastpage}
\usepackage{tasks}
\usepackage{needspace}
\usepackage[most]{tcolorbox}
\usepackage{array}
\usepackage[hidelinks,pdftitle={Fitxa d'exercicis de la unitat 7: Funcions elementals. 1r de Batxillerat MI},pdfsubject={Matemàtiques I},pdfkeywords={funció lineal, funció afí, paràbola, funció cúbica, hipèrbola, funció radical, exponencial, logaritme, funcions trigonomètriques, funcions a trossos, valor absolut, problemes d'optimització, ingressos i costos}]{hyperref}

\definecolor{navy}{HTML}{1B2748}
\definecolor{accent}{HTML}{E8603A}
\definecolor{teal}{HTML}{1F9E90}
\definecolor{soft}{HTML}{EEF1F8}
\definecolor{muted}{HTML}{5D6685}

\setlength{\parindent}{0pt}
\setlength{\parskip}{3pt}
\settasks{label=\alph*), label-format=\bfseries\color{navy}, label-width=1.3em, item-indent=1.8em, column-sep=0.6em, after-item-skip=5pt, before-skip=3pt, after-skip=4pt}

% ---------- capçalera i peu ----------
\newcommand{\llicencia}{\href{https://creativecommons.org/licenses/by-nc-sa/4.0/deed.ca}{\ccbyncsa}\ \ Llicència Creative Commons Reconeixement-NoComercial-CompartirIgual 4.0 Internacional (CC BY-NC-SA 4.0)}
\pagestyle{fancy}
\fancyhf{}
\fancyhead[L]{\raisebox{-0.35\height}{\includegraphics[height=0.95cm]{logo-gva.png}}}
\fancyhead[C]{\small\color{muted}Fitxa d'exercicis · Unitat 7: Funcions elementals · 1r Batxillerat MI}
\fancyhead[R]{\raisebox{-0.35\height}{\includegraphics[height=0.95cm]{logo-vermellar.png}}}
\fancyfoot[L]{\scriptsize\color{muted}\llicencia}
\fancyfoot[R]{\small\color{muted}\thepage}
\renewcommand{\headrulewidth}{0.4pt}
\renewcommand{\headrule}{\color{muted!40}\hrule height \headrulewidth}
\fancypagestyle{primera}{\fancyhf{}\renewcommand{\headrulewidth}{0pt}
  \fancyfoot[L]{\scriptsize\color{muted}\llicencia}
  \fancyfoot[R]{\small\color{muted}\thepage}}

% ---------- nivells i exercicis ----------
\newcommand{\estrelles}[1]{{\color{accent}\ifcase#1\or\faStar\faStar[regular]\faStar[regular]\or\faStar\faStar\faStar[regular]\or\faStar\faStar\faStar\fi}}
\newcounter{exer}
\newcommand{\ex}[1]{\par\needspace{4\baselineskip}\medskip\refstepcounter{exer}\noindent{\color{navy}\bfseries\theexer.}\ \ }

\newcommand{\seccio}[2]{\par\needspace{16\baselineskip}\bigskip
  \begin{tcolorbox}[colback=navy,colframe=navy,coltext=white,arc=2mm,boxrule=0pt,left=3mm,right=3mm,top=1.2mm,bottom=1.2mm]
  {\large\bfseries #1\quad #2}
  \end{tcolorbox}\vspace{-2pt}}

\newcommand{\exf}[1]{\par\medskip\refstepcounter{exer}\noindent{\color{navy}\bfseries\theexer.}\ \ }
\usetikzlibrary{calc}
\newcolumntype{C}[1]{>{\centering\arraybackslash}p{#1}}
\providecommand{\gt}{>}\providecommand{\lt}{<}
% intervals a la manera dels apunts: ]a,b[, [a,b], ...
\newcommand{\R}{\mathbb{R}}
\newcommand{\oo}[2]{\left]#1,\,#2\right[}
\newcommand{\oc}[2]{\left]#1,\,#2\right]}
\newcommand{\co}[2]{\left[#1,\,#2\right[}
\newcommand{\cc}[2]{\left[#1,\,#2\right]}
% quadre de recordatori (només a la fitxa d'exercicis)
\newcommand{\recorda}[1]{\par\smallskip\begin{tcolorbox}[colback=accent!7,colframe=accent!55,arc=2mm,boxrule=0.5pt,left=2.5mm,right=2.5mm,top=1mm,bottom=1mm]\small{\color{accent}\faLightbulb[regular]}\ \textbf{Recorda:}\ #1\end{tcolorbox}}
\newcommand{\res}[1]{#1}\newcommand{\rest}[1]{#1}

\def\fmneed{11}% que el títol d'una secció no quede sol al final del full
% =====================================================================
% Capçalera FentMates de les fitxes: un color per a cada curs.
% Els colors són els de la pàgina inicial (LEVELS a assets/js/materials.js);
% el logo ve de _eines/logo/pdf/ (python3 _eines/logo/genera.py).
%
% Ús, al final del preàmbul de la unitat:
%   % =====================================================================
% Capçalera FentMates de les fitxes: un color per a cada curs.
% Els colors són els de la pàgina inicial (LEVELS a assets/js/materials.js);
% el logo ve de _eines/logo/pdf/ (python3 _eines/logo/genera.py).
%
% Ús, al final del preàmbul de la unitat:
%   % =====================================================================
% Capçalera FentMates de les fitxes: un color per a cada curs.
% Els colors són els de la pàgina inicial (LEVELS a assets/js/materials.js);
% el logo ve de _eines/logo/pdf/ (python3 _eines/logo/genera.py).
%
% Ús, al final del preàmbul de la unitat:
%   \input{../comu/fentmates.tex}
%   \fmcurs{1eso}
% i, al començament del document:
%   \capcalera{Matemàtiques}{1r ESO}{Fitxa d'exercicis · Unitat 2}{Divisibilitat}
% Opcional abans de \input: \def\fmneed{14} (espai mínim abans de cada secció).
% =====================================================================
\makeatletter
\def\fm@logodir{../../logo/pdf}
\providecommand{\fmneed}{10}
\def\fm@cursnom{}

% clau del curs, color i color del text damunt del color
\newcommand{\fm@curs}[3]{\@namedef{fm@color@#1}{#2}\@namedef{fm@text@#1}{#3}}
\fm@curs{1eso}{E6B84A}{navy}
\fm@curs{2eso}{3B82F6}{white}
\fm@curs{3eso}{FF7A50}{white}
\fm@curs{4eso-a}{A78BFA}{white}
\fm@curs{4eso-b}{6D28D9}{white}
\fm@curs{1batx-mi}{9B1C1C}{white}
\fm@curs{1batx-ccss}{EF6B6B}{white}
\fm@curs{2batx-mii}{1E6B47}{white}
\fm@curs{2batx-ccss}{2BB3A3}{white}

\newcommand{\fmcurs}[1]{%
  \@ifundefined{fm@color@#1}{\PackageError{fentmates}{Curs desconegut: #1}{}}{}%
  \edef\fm@tmp{\noexpand\definecolor{curs}{HTML}{\@nameuse{fm@color@#1}}%
               \noexpand\colorlet{curstext}{\@nameuse{fm@text@#1}}}\fm@tmp
  \colorlet{cursfosc}{curs!72!black}%
  \def\fm@logo{\fm@logodir/vertical-#1.pdf}%
  \def\fm@marca{\fm@logodir/marca-#1.pdf}}

% ---------- capçalera de la primera pàgina ----------
% #1 matèria (petit), #2 curs, #3 tipus de document i unitat, #4 títol
\newcommand{\capcalera}[4]{%
  \gdef\fm@cursnom{#2}%
  \par\noindent
  \begin{tcolorbox}[colback=curs,colframe=curs,boxrule=0pt,arc=2.2mm,width=3.5cm,height=2.15cm,
      valign=center,left=3mm,right=2mm,top=1mm,bottom=1mm,nobeforeafter,box align=center]
    \color{curstext}{\footnotesize #1}\par\vspace{1pt}{\fontsize{19}{21}\selectfont\bfseries #2\par}%
  \end{tcolorbox}\hspace{5mm}%
  \begin{minipage}[c]{\dimexpr\linewidth-3.5cm-5mm-2.9cm\relax}
    {\color{cursfosc}\small\bfseries\MakeUppercase{#3}}\par\vspace{2pt}
    {\color{navy}\fontsize{20}{23}\selectfont\bfseries #4\par}
  \end{minipage}\hfill
  \begin{minipage}[c]{2.6cm}\raggedleft\includegraphics[width=2.6cm]{\fm@logo}\end{minipage}\par
  \vspace{2pt}}

% ---------- barres de secció del color del curs ----------
\renewcommand{\seccio}[2]{\par\needspace{\fmneed\baselineskip}\bigskip
  \begin{tcolorbox}[colback=curs,colframe=curs,coltext=curstext,arc=2mm,boxrule=0pt,left=3mm,right=3mm,top=1.2mm,bottom=1.2mm]
  {\large\bfseries #1\quad #2}
  \end{tcolorbox}\vspace{-2pt}}

% ---------- estrelles de dificultat del color del curs (a Batxillerat no se'n posen) ----------
\@ifundefined{estrelles}{}{%
  \renewcommand{\estrelles}[1]{{\color{curs}\ifcase#1\or\faStar\faStar[regular]\faStar[regular]\or\faStar\faStar\faStar[regular]\or\faStar\faStar\faStar\fi}}}

% ---------- capçalera i peu de les pàgines ----------
% primera pàgina: logos institucionals al peu
\fancypagestyle{primera}{\fancyhf{}\renewcommand{\headrulewidth}{0pt}
  \fancyfoot[L]{\raisebox{-0.32\height}{\includegraphics[height=0.78cm]{logo-gva.png}}\hspace{5mm}%
    \raisebox{-0.32\height}{\includegraphics[height=0.78cm]{logo-vermellar.png}}}
  \fancyfoot[R]{\scriptsize\color{muted}\href{https://creativecommons.org/licenses/by-nc-sa/4.0/deed.ca}{\ccbyncsa}\ \ CC BY-NC-SA 4.0 · fentmates.org\qquad\small\thepage}}
% la resta: el símbol a l'esquerra i el curs a la dreta, amb el títol al centre com abans
\fancyhead[L]{\raisebox{-0.3\height}{\includegraphics[height=0.62cm]{\fm@marca}}}
\fancyhead[R]{\small\bfseries\color{cursfosc}\fm@cursnom}
\renewcommand{\headrule}{\color{curs!45}\hrule height \headrulewidth}
\makeatother

%   \fmcurs{1eso}
% i, al començament del document:
%   \capcalera{Matemàtiques}{1r ESO}{Fitxa d'exercicis · Unitat 2}{Divisibilitat}
% Opcional abans de \input: \def\fmneed{14} (espai mínim abans de cada secció).
% =====================================================================
\makeatletter
\def\fm@logodir{../../logo/pdf}
\providecommand{\fmneed}{10}
\def\fm@cursnom{}

% clau del curs, color i color del text damunt del color
\newcommand{\fm@curs}[3]{\@namedef{fm@color@#1}{#2}\@namedef{fm@text@#1}{#3}}
\fm@curs{1eso}{E6B84A}{navy}
\fm@curs{2eso}{3B82F6}{white}
\fm@curs{3eso}{FF7A50}{white}
\fm@curs{4eso-a}{A78BFA}{white}
\fm@curs{4eso-b}{6D28D9}{white}
\fm@curs{1batx-mi}{9B1C1C}{white}
\fm@curs{1batx-ccss}{EF6B6B}{white}
\fm@curs{2batx-mii}{1E6B47}{white}
\fm@curs{2batx-ccss}{2BB3A3}{white}

\newcommand{\fmcurs}[1]{%
  \@ifundefined{fm@color@#1}{\PackageError{fentmates}{Curs desconegut: #1}{}}{}%
  \edef\fm@tmp{\noexpand\definecolor{curs}{HTML}{\@nameuse{fm@color@#1}}%
               \noexpand\colorlet{curstext}{\@nameuse{fm@text@#1}}}\fm@tmp
  \colorlet{cursfosc}{curs!72!black}%
  \def\fm@logo{\fm@logodir/vertical-#1.pdf}%
  \def\fm@marca{\fm@logodir/marca-#1.pdf}}

% ---------- capçalera de la primera pàgina ----------
% #1 matèria (petit), #2 curs, #3 tipus de document i unitat, #4 títol
\newcommand{\capcalera}[4]{%
  \gdef\fm@cursnom{#2}%
  \par\noindent
  \begin{tcolorbox}[colback=curs,colframe=curs,boxrule=0pt,arc=2.2mm,width=3.5cm,height=2.15cm,
      valign=center,left=3mm,right=2mm,top=1mm,bottom=1mm,nobeforeafter,box align=center]
    \color{curstext}{\footnotesize #1}\par\vspace{1pt}{\fontsize{19}{21}\selectfont\bfseries #2\par}%
  \end{tcolorbox}\hspace{5mm}%
  \begin{minipage}[c]{\dimexpr\linewidth-3.5cm-5mm-2.9cm\relax}
    {\color{cursfosc}\small\bfseries\MakeUppercase{#3}}\par\vspace{2pt}
    {\color{navy}\fontsize{20}{23}\selectfont\bfseries #4\par}
  \end{minipage}\hfill
  \begin{minipage}[c]{2.6cm}\raggedleft\includegraphics[width=2.6cm]{\fm@logo}\end{minipage}\par
  \vspace{2pt}}

% ---------- barres de secció del color del curs ----------
\renewcommand{\seccio}[2]{\par\needspace{\fmneed\baselineskip}\bigskip
  \begin{tcolorbox}[colback=curs,colframe=curs,coltext=curstext,arc=2mm,boxrule=0pt,left=3mm,right=3mm,top=1.2mm,bottom=1.2mm]
  {\large\bfseries #1\quad #2}
  \end{tcolorbox}\vspace{-2pt}}

% ---------- estrelles de dificultat del color del curs (a Batxillerat no se'n posen) ----------
\@ifundefined{estrelles}{}{%
  \renewcommand{\estrelles}[1]{{\color{curs}\ifcase#1\or\faStar\faStar[regular]\faStar[regular]\or\faStar\faStar\faStar[regular]\or\faStar\faStar\faStar\fi}}}

% ---------- capçalera i peu de les pàgines ----------
% primera pàgina: logos institucionals al peu
\fancypagestyle{primera}{\fancyhf{}\renewcommand{\headrulewidth}{0pt}
  \fancyfoot[L]{\raisebox{-0.32\height}{\includegraphics[height=0.78cm]{logo-gva.png}}\hspace{5mm}%
    \raisebox{-0.32\height}{\includegraphics[height=0.78cm]{logo-vermellar.png}}}
  \fancyfoot[R]{\scriptsize\color{muted}\href{https://creativecommons.org/licenses/by-nc-sa/4.0/deed.ca}{\ccbyncsa}\ \ CC BY-NC-SA 4.0 · fentmates.org\qquad\small\thepage}}
% la resta: el símbol a l'esquerra i el curs a la dreta, amb el títol al centre com abans
\fancyhead[L]{\raisebox{-0.3\height}{\includegraphics[height=0.62cm]{\fm@marca}}}
\fancyhead[R]{\small\bfseries\color{cursfosc}\fm@cursnom}
\renewcommand{\headrule}{\color{curs!45}\hrule height \headrulewidth}
\makeatother

%   \fmcurs{1eso}
% i, al començament del document:
%   \capcalera{Matemàtiques}{1r ESO}{Fitxa d'exercicis · Unitat 2}{Divisibilitat}
% Opcional abans de \input: \def\fmneed{14} (espai mínim abans de cada secció).
% =====================================================================
\makeatletter
\def\fm@logodir{../../logo/pdf}
\providecommand{\fmneed}{10}
\def\fm@cursnom{}

% clau del curs, color i color del text damunt del color
\newcommand{\fm@curs}[3]{\@namedef{fm@color@#1}{#2}\@namedef{fm@text@#1}{#3}}
\fm@curs{1eso}{E6B84A}{navy}
\fm@curs{2eso}{3B82F6}{white}
\fm@curs{3eso}{FF7A50}{white}
\fm@curs{4eso-a}{A78BFA}{white}
\fm@curs{4eso-b}{6D28D9}{white}
\fm@curs{1batx-mi}{9B1C1C}{white}
\fm@curs{1batx-ccss}{EF6B6B}{white}
\fm@curs{2batx-mii}{1E6B47}{white}
\fm@curs{2batx-ccss}{2BB3A3}{white}

\newcommand{\fmcurs}[1]{%
  \@ifundefined{fm@color@#1}{\PackageError{fentmates}{Curs desconegut: #1}{}}{}%
  \edef\fm@tmp{\noexpand\definecolor{curs}{HTML}{\@nameuse{fm@color@#1}}%
               \noexpand\colorlet{curstext}{\@nameuse{fm@text@#1}}}\fm@tmp
  \colorlet{cursfosc}{curs!72!black}%
  \def\fm@logo{\fm@logodir/vertical-#1.pdf}%
  \def\fm@marca{\fm@logodir/marca-#1.pdf}}

% ---------- capçalera de la primera pàgina ----------
% #1 matèria (petit), #2 curs, #3 tipus de document i unitat, #4 títol
\newcommand{\capcalera}[4]{%
  \gdef\fm@cursnom{#2}%
  \par\noindent
  \begin{tcolorbox}[colback=curs,colframe=curs,boxrule=0pt,arc=2.2mm,width=3.5cm,height=2.15cm,
      valign=center,left=3mm,right=2mm,top=1mm,bottom=1mm,nobeforeafter,box align=center]
    \color{curstext}{\footnotesize #1}\par\vspace{1pt}{\fontsize{19}{21}\selectfont\bfseries #2\par}%
  \end{tcolorbox}\hspace{5mm}%
  \begin{minipage}[c]{\dimexpr\linewidth-3.5cm-5mm-2.9cm\relax}
    {\color{cursfosc}\small\bfseries\MakeUppercase{#3}}\par\vspace{2pt}
    {\color{navy}\fontsize{20}{23}\selectfont\bfseries #4\par}
  \end{minipage}\hfill
  \begin{minipage}[c]{2.6cm}\raggedleft\includegraphics[width=2.6cm]{\fm@logo}\end{minipage}\par
  \vspace{2pt}}

% ---------- barres de secció del color del curs ----------
\renewcommand{\seccio}[2]{\par\needspace{\fmneed\baselineskip}\bigskip
  \begin{tcolorbox}[colback=curs,colframe=curs,coltext=curstext,arc=2mm,boxrule=0pt,left=3mm,right=3mm,top=1.2mm,bottom=1.2mm]
  {\large\bfseries #1\quad #2}
  \end{tcolorbox}\vspace{-2pt}}

% ---------- estrelles de dificultat del color del curs (a Batxillerat no se'n posen) ----------
\@ifundefined{estrelles}{}{%
  \renewcommand{\estrelles}[1]{{\color{curs}\ifcase#1\or\faStar\faStar[regular]\faStar[regular]\or\faStar\faStar\faStar[regular]\or\faStar\faStar\faStar\fi}}}

% ---------- capçalera i peu de les pàgines ----------
% primera pàgina: logos institucionals al peu
\fancypagestyle{primera}{\fancyhf{}\renewcommand{\headrulewidth}{0pt}
  \fancyfoot[L]{\raisebox{-0.32\height}{\includegraphics[height=0.78cm]{logo-gva.png}}\hspace{5mm}%
    \raisebox{-0.32\height}{\includegraphics[height=0.78cm]{logo-vermellar.png}}}
  \fancyfoot[R]{\scriptsize\color{muted}\href{https://creativecommons.org/licenses/by-nc-sa/4.0/deed.ca}{\ccbyncsa}\ \ CC BY-NC-SA 4.0 · fentmates.org\qquad\small\thepage}}
% la resta: el símbol a l'esquerra i el curs a la dreta, amb el títol al centre com abans
\fancyhead[L]{\raisebox{-0.3\height}{\includegraphics[height=0.62cm]{\fm@marca}}}
\fancyhead[R]{\small\bfseries\color{cursfosc}\fm@cursnom}
\renewcommand{\headrule}{\color{curs!45}\hrule height \headrulewidth}
\makeatother

\fmcurs{1batx-mi}

\begin{document}
\thispagestyle{primera}

\capcalera{Matemàtiques I}{1r Batx}{Fitxa d'exercicis · Unitat 7}{\raggedright Funcions elementals}
\vspace{10pt}
Nom i cognoms: \rule{8.4cm}{0.4pt}\hfill Grup: \rule{1.6cm}{0.4pt}\hfill Data: \rule{2.2cm}{0.4pt}
\vspace{8pt}

\begin{tcolorbox}[colback=soft,colframe=soft,arc=2mm,left=3mm,right=3mm,top=1.5mm,bottom=1.5mm]
\small Resol els exercicis al quadern mostrant tots els passos i fes un dibuix sempre que t'ajude. Escriu els intervals com als apunts ($\left]a,b\right[$ obert, $[a,b]$ tancat). Les seccions segueixen l'ordre dels apunts i la darrera té problemes de màxims i mínims amb funcions quadràtiques.
\end{tcolorbox}

% =====================================================
\seccio{1}{Funcions lineals i afins}

\exf{1} Troba l'equació de la recta que passa per $A(-2,1)$ i $B(4,4)$. És creixent o decreixent?

\ex{1} Troba l'equació de la recta de pendent $-3$ que passa per $(2,\,-1)$ i els seus punts de tall amb els eixos.

\ex{1} Escriu l'equació de cadascuna d'aquestes rectes:\begin{tasks}(3)
  \task \begin{tikzpicture}[x=0.486cm,y=0.486cm,baseline={(current bounding box.center)},font=\scriptsize]
\draw[muted!25,line width=.3pt] (-3,-4)--(-3,3);
\draw[muted!25,line width=.3pt] (-2,-4)--(-2,3);
\draw[muted!25,line width=.3pt] (-1,-4)--(-1,3);
\draw[muted!25,line width=.3pt] (0,-4)--(0,3);
\draw[muted!25,line width=.3pt] (1,-4)--(1,3);
\draw[muted!25,line width=.3pt] (2,-4)--(2,3);
\draw[muted!25,line width=.3pt] (3,-4)--(3,3);
\draw[muted!25,line width=.3pt] (4,-4)--(4,3);
\draw[muted!25,line width=.3pt] (-3,-4)--(4,-4);
\draw[muted!25,line width=.3pt] (-3,-3)--(4,-3);
\draw[muted!25,line width=.3pt] (-3,-2)--(4,-2);
\draw[muted!25,line width=.3pt] (-3,-1)--(4,-1);
\draw[muted!25,line width=.3pt] (-3,0)--(4,0);
\draw[muted!25,line width=.3pt] (-3,1)--(4,1);
\draw[muted!25,line width=.3pt] (-3,2)--(4,2);
\draw[muted!25,line width=.3pt] (-3,3)--(4,3);
\draw[-latex,muted] (-3,0)--(4,0);
\draw[-latex,muted] (0,-4)--(0,3);
\node[below=-1pt,text=muted] at (-3,0) {$-3$};
\node[below=-1pt,text=muted] at (-2,0) {$-2$};
\node[below=-1pt,text=muted] at (-1,0) {$-1$};
\node[below=-1pt,text=muted] at (1,0) {$1$};
\node[below=-1pt,text=muted] at (2,0) {$2$};
\node[below=-1pt,text=muted] at (3,0) {$3$};
\node[below=-1pt,text=muted] at (4,0) {$4$};
\node[left=-1pt,text=muted] at (0,-4) {$-4$};
\node[left=-1pt,text=muted] at (0,-3) {$-3$};
\node[left=-1pt,text=muted] at (0,-2) {$-2$};
\node[left=-1pt,text=muted] at (0,-1) {$-1$};
\node[left=-1pt,text=muted] at (0,1) {$1$};
\node[left=-1pt,text=muted] at (0,2) {$2$};
\node[left=-1pt,text=muted] at (0,3) {$3$};
\begin{scope}\clip (-3,-4) rectangle (4,3);
\draw[accent,line width=1.1pt,line join=round] (-3,-9)--(-2.944,-8.888)--(-2.888,-8.776)--(-2.832,-8.664)--(-2.776,-8.552)--(-2.72,-8.44)--(-2.664,-8.328)--(-2.608,-8.216)--(-2.552,-8.104)--(-2.496,-7.992)--(-2.44,-7.88)--(-2.384,-7.768)--(-2.328,-7.656)--(-2.272,-7.544)--(-2.216,-7.432)--(-2.16,-7.32)--(-2.104,-7.208)--(-2.048,-7.096)--(-1.992,-6.984)--(-1.936,-6.872)--(-1.88,-6.76)--(-1.824,-6.648)--(-1.768,-6.536)--(-1.712,-6.424)--(-1.656,-6.312)--(-1.6,-6.2)--(-1.544,-6.088)--(-1.488,-5.976)--(-1.432,-5.864)--(-1.376,-5.752)--(-1.32,-5.64)--(-1.264,-5.528)--(-1.208,-5.416)--(-1.152,-5.304)--(-1.096,-5.192)--(-1.04,-5.08)--(-0.984,-4.968)--(-0.928,-4.856)--(-0.872,-4.744)--(-0.816,-4.632)--(-0.76,-4.52)--(-0.704,-4.408)--(-0.648,-4.296)--(-0.592,-4.184)--(-0.536,-4.072)--(-0.48,-3.96)--(-0.424,-3.848)--(-0.368,-3.736)--(-0.312,-3.624)--(-0.256,-3.512)--(-0.2,-3.4)--(-0.144,-3.288)--(-0.088,-3.176)--(-0.032,-3.064)--(0.024,-2.952)--(0.08,-2.84)--(0.136,-2.728)--(0.192,-2.616)--(0.248,-2.504)--(0.304,-2.392)--(0.36,-2.28)--(0.416,-2.168)--(0.472,-2.056)--(0.528,-1.944)--(0.584,-1.832)--(0.64,-1.72)--(0.696,-1.608)--(0.752,-1.496)--(0.808,-1.384)--(0.864,-1.272)--(0.92,-1.16)--(0.976,-1.048)--(1.032,-0.936)--(1.088,-0.824)--(1.144,-0.712)--(1.2,-0.6)--(1.256,-0.488)--(1.312,-0.376)--(1.368,-0.264)--(1.424,-0.152)--(1.48,-0.04)--(1.536,0.072)--(1.592,0.184)--(1.648,0.296)--(1.704,0.408)--(1.76,0.52)--(1.816,0.632)--(1.872,0.744)--(1.928,0.856)--(1.984,0.968)--(2.04,1.08)--(2.096,1.192)--(2.152,1.304)--(2.208,1.416)--(2.264,1.528)--(2.32,1.64)--(2.376,1.752)--(2.432,1.864)--(2.488,1.976)--(2.544,2.088)--(2.6,2.2)--(2.656,2.312)--(2.712,2.424)--(2.768,2.536)--(2.824,2.648)--(2.88,2.76)--(2.936,2.872)--(2.992,2.984)--(3.048,3.096)--(3.104,3.208)--(3.16,3.32)--(3.216,3.432)--(3.272,3.544)--(3.328,3.656)--(3.384,3.768)--(3.44,3.88)--(3.496,3.992)--(3.552,4.104)--(3.608,4.216)--(3.664,4.328)--(3.72,4.44)--(3.776,4.552)--(3.832,4.664)--(3.888,4.776)--(3.944,4.888)--(4,5);
\filldraw[fill=blue!70!black,draw=blue!70!black,line width=.8pt] (0,-3) circle (1.5pt);
\filldraw[fill=blue!70!black,draw=blue!70!black,line width=.8pt] (2,1) circle (1.5pt);
\end{scope}\end{tikzpicture}
  \task \begin{tikzpicture}[x=0.378cm,y=0.422cm,baseline={(current bounding box.center)},font=\scriptsize]
\draw[muted!25,line width=.3pt] (-2,-1)--(-2,4);
\draw[muted!25,line width=.3pt] (-1,-1)--(-1,4);
\draw[muted!25,line width=.3pt] (0,-1)--(0,4);
\draw[muted!25,line width=.3pt] (1,-1)--(1,4);
\draw[muted!25,line width=.3pt] (2,-1)--(2,4);
\draw[muted!25,line width=.3pt] (3,-1)--(3,4);
\draw[muted!25,line width=.3pt] (4,-1)--(4,4);
\draw[muted!25,line width=.3pt] (5,-1)--(5,4);
\draw[muted!25,line width=.3pt] (6,-1)--(6,4);
\draw[muted!25,line width=.3pt] (7,-1)--(7,4);
\draw[muted!25,line width=.3pt] (-2,-1)--(7,-1);
\draw[muted!25,line width=.3pt] (-2,0)--(7,0);
\draw[muted!25,line width=.3pt] (-2,1)--(7,1);
\draw[muted!25,line width=.3pt] (-2,2)--(7,2);
\draw[muted!25,line width=.3pt] (-2,3)--(7,3);
\draw[muted!25,line width=.3pt] (-2,4)--(7,4);
\draw[-latex,muted] (-2,0)--(7,0);
\draw[-latex,muted] (0,-1)--(0,4);
\node[below=-1pt,text=muted] at (-2,0) {$-2$};
\node[below=-1pt,text=muted] at (-1,0) {$-1$};
\node[below=-1pt,text=muted] at (1,0) {$1$};
\node[below=-1pt,text=muted] at (2,0) {$2$};
\node[below=-1pt,text=muted] at (3,0) {$3$};
\node[below=-1pt,text=muted] at (4,0) {$4$};
\node[below=-1pt,text=muted] at (5,0) {$5$};
\node[below=-1pt,text=muted] at (6,0) {$6$};
\node[below=-1pt,text=muted] at (7,0) {$7$};
\node[left=-1pt,text=muted] at (0,-1) {$-1$};
\node[left=-1pt,text=muted] at (0,1) {$1$};
\node[left=-1pt,text=muted] at (0,2) {$2$};
\node[left=-1pt,text=muted] at (0,3) {$3$};
\node[left=-1pt,text=muted] at (0,4) {$4$};
\begin{scope}\clip (-2,-1) rectangle (7,4);
\draw[accent,line width=1.1pt,line join=round] (-2,2.667)--(-1.928,2.643)--(-1.856,2.619)--(-1.784,2.595)--(-1.712,2.571)--(-1.64,2.547)--(-1.568,2.523)--(-1.496,2.499)--(-1.424,2.475)--(-1.352,2.451)--(-1.28,2.427)--(-1.208,2.403)--(-1.136,2.379)--(-1.064,2.355)--(-0.992,2.331)--(-0.92,2.307)--(-0.848,2.283)--(-0.776,2.259)--(-0.704,2.235)--(-0.632,2.211)--(-0.56,2.187)--(-0.488,2.163)--(-0.416,2.139)--(-0.344,2.115)--(-0.272,2.091)--(-0.2,2.067)--(-0.128,2.043)--(-0.056,2.019)--(0.016,1.995)--(0.088,1.971)--(0.16,1.947)--(0.232,1.923)--(0.304,1.899)--(0.376,1.875)--(0.448,1.851)--(0.52,1.827)--(0.592,1.803)--(0.664,1.779)--(0.736,1.755)--(0.808,1.731)--(0.88,1.707)--(0.952,1.683)--(1.024,1.659)--(1.096,1.635)--(1.168,1.611)--(1.24,1.587)--(1.312,1.563)--(1.384,1.539)--(1.456,1.515)--(1.528,1.491)--(1.6,1.467)--(1.672,1.443)--(1.744,1.419)--(1.816,1.395)--(1.888,1.371)--(1.96,1.347)--(2.032,1.323)--(2.104,1.299)--(2.176,1.275)--(2.248,1.251)--(2.32,1.227)--(2.392,1.203)--(2.464,1.179)--(2.536,1.155)--(2.608,1.131)--(2.68,1.107)--(2.752,1.083)--(2.824,1.059)--(2.896,1.035)--(2.968,1.011)--(3.04,0.987)--(3.112,0.963)--(3.184,0.939)--(3.256,0.915)--(3.328,0.891)--(3.4,0.867)--(3.472,0.843)--(3.544,0.819)--(3.616,0.795)--(3.688,0.771)--(3.76,0.747)--(3.832,0.723)--(3.904,0.699)--(3.976,0.675)--(4.048,0.651)--(4.12,0.627)--(4.192,0.603)--(4.264,0.579)--(4.336,0.555)--(4.408,0.531)--(4.48,0.507)--(4.552,0.483)--(4.624,0.459)--(4.696,0.435)--(4.768,0.411)--(4.84,0.387)--(4.912,0.363)--(4.984,0.339)--(5.056,0.315)--(5.128,0.291)--(5.2,0.267)--(5.272,0.243)--(5.344,0.219)--(5.416,0.195)--(5.488,0.171)--(5.56,0.147)--(5.632,0.123)--(5.704,0.099)--(5.776,0.075)--(5.848,0.051)--(5.92,0.027)--(5.992,0.003)--(6.064,-0.021)--(6.136,-0.045)--(6.208,-0.069)--(6.28,-0.093)--(6.352,-0.117)--(6.424,-0.141)--(6.496,-0.165)--(6.568,-0.189)--(6.64,-0.213)--(6.712,-0.237)--(6.784,-0.261)--(6.856,-0.285)--(6.928,-0.309)--(7,-0.333);
\filldraw[fill=blue!70!black,draw=blue!70!black,line width=.8pt] (0,2) circle (1.5pt);
\filldraw[fill=blue!70!black,draw=blue!70!black,line width=.8pt] (3,1) circle (1.5pt);
\end{scope}\end{tikzpicture}
  \task \begin{tikzpicture}[x=0.425cm,y=0.425cm,baseline={(current bounding box.center)},font=\scriptsize]
\draw[muted!25,line width=.3pt] (-4,-3)--(-4,2);
\draw[muted!25,line width=.3pt] (-3,-3)--(-3,2);
\draw[muted!25,line width=.3pt] (-2,-3)--(-2,2);
\draw[muted!25,line width=.3pt] (-1,-3)--(-1,2);
\draw[muted!25,line width=.3pt] (0,-3)--(0,2);
\draw[muted!25,line width=.3pt] (1,-3)--(1,2);
\draw[muted!25,line width=.3pt] (2,-3)--(2,2);
\draw[muted!25,line width=.3pt] (3,-3)--(3,2);
\draw[muted!25,line width=.3pt] (4,-3)--(4,2);
\draw[muted!25,line width=.3pt] (-4,-3)--(4,-3);
\draw[muted!25,line width=.3pt] (-4,-2)--(4,-2);
\draw[muted!25,line width=.3pt] (-4,-1)--(4,-1);
\draw[muted!25,line width=.3pt] (-4,0)--(4,0);
\draw[muted!25,line width=.3pt] (-4,1)--(4,1);
\draw[muted!25,line width=.3pt] (-4,2)--(4,2);
\draw[-latex,muted] (-4,0)--(4,0);
\draw[-latex,muted] (0,-3)--(0,2);
\node[below=-1pt,text=muted] at (-4,0) {$-4$};
\node[below=-1pt,text=muted] at (-3,0) {$-3$};
\node[below=-1pt,text=muted] at (-2,0) {$-2$};
\node[below=-1pt,text=muted] at (-1,0) {$-1$};
\node[below=-1pt,text=muted] at (1,0) {$1$};
\node[below=-1pt,text=muted] at (2,0) {$2$};
\node[below=-1pt,text=muted] at (3,0) {$3$};
\node[below=-1pt,text=muted] at (4,0) {$4$};
\node[left=-1pt,text=muted] at (0,-3) {$-3$};
\node[left=-1pt,text=muted] at (0,-2) {$-2$};
\node[left=-1pt,text=muted] at (0,-1) {$-1$};
\node[left=-1pt,text=muted] at (0,1) {$1$};
\node[left=-1pt,text=muted] at (0,2) {$2$};
\begin{scope}\clip (-4,-3) rectangle (4,2);
\draw[accent,line width=1.1pt,line join=round] (-4,-2)--(-3.936,-2)--(-3.872,-2)--(-3.808,-2)--(-3.744,-2)--(-3.68,-2)--(-3.616,-2)--(-3.552,-2)--(-3.488,-2)--(-3.424,-2)--(-3.36,-2)--(-3.296,-2)--(-3.232,-2)--(-3.168,-2)--(-3.104,-2)--(-3.04,-2)--(-2.976,-2)--(-2.912,-2)--(-2.848,-2)--(-2.784,-2)--(-2.72,-2)--(-2.656,-2)--(-2.592,-2)--(-2.528,-2)--(-2.464,-2)--(-2.4,-2)--(-2.336,-2)--(-2.272,-2)--(-2.208,-2)--(-2.144,-2)--(-2.08,-2)--(-2.016,-2)--(-1.952,-2)--(-1.888,-2)--(-1.824,-2)--(-1.76,-2)--(-1.696,-2)--(-1.632,-2)--(-1.568,-2)--(-1.504,-2)--(-1.44,-2)--(-1.376,-2)--(-1.312,-2)--(-1.248,-2)--(-1.184,-2)--(-1.12,-2)--(-1.056,-2)--(-0.992,-2)--(-0.928,-2)--(-0.864,-2)--(-0.8,-2)--(-0.736,-2)--(-0.672,-2)--(-0.608,-2)--(-0.544,-2)--(-0.48,-2)--(-0.416,-2)--(-0.352,-2)--(-0.288,-2)--(-0.224,-2)--(-0.16,-2)--(-0.096,-2)--(-0.032,-2)--(0.032,-2)--(0.096,-2)--(0.16,-2)--(0.224,-2)--(0.288,-2)--(0.352,-2)--(0.416,-2)--(0.48,-2)--(0.544,-2)--(0.608,-2)--(0.672,-2)--(0.736,-2)--(0.8,-2)--(0.864,-2)--(0.928,-2)--(0.992,-2)--(1.056,-2)--(1.12,-2)--(1.184,-2)--(1.248,-2)--(1.312,-2)--(1.376,-2)--(1.44,-2)--(1.504,-2)--(1.568,-2)--(1.632,-2)--(1.696,-2)--(1.76,-2)--(1.824,-2)--(1.888,-2)--(1.952,-2)--(2.016,-2)--(2.08,-2)--(2.144,-2)--(2.208,-2)--(2.272,-2)--(2.336,-2)--(2.4,-2)--(2.464,-2)--(2.528,-2)--(2.592,-2)--(2.656,-2)--(2.72,-2)--(2.784,-2)--(2.848,-2)--(2.912,-2)--(2.976,-2)--(3.04,-2)--(3.104,-2)--(3.168,-2)--(3.232,-2)--(3.296,-2)--(3.36,-2)--(3.424,-2)--(3.488,-2)--(3.552,-2)--(3.616,-2)--(3.68,-2)--(3.744,-2)--(3.808,-2)--(3.872,-2)--(3.936,-2)--(4,-2);
\end{scope}\end{tikzpicture}
\end{tasks}

\ex{2} Un taxi cobra $3{,}50$~€ de baixada de bandera i $1{,}20$~€ per quilòmetre. Escriu la funció que dona el preu en funció dels quilòmetres. Quant costa un viatge de 12 km? Quants quilòmetres has fet si has pagat 23~€?

\ex{1} Troba l'equació de la recta paral·lela a $y=3x-1$ que passa pel punt $(1,5)$.

\ex{2} Un dipòsit té 200 L d'aigua i se'n buida a raó de 8 L per minut. Escriu la funció $V(t)$ que dona els litres que queden al cap de $t$ minuts. Quan es buida? Quin és el domini de la funció en aquest context?


% =====================================================
\seccio{2}{Funcions quadràtiques i cúbiques}

\exf{1} Representa aquestes paràboles calculant-ne la forma, el vèrtex i els punts de tall:\begin{tasks}(4)
  \task $y=x^2-6x+5$
  \task $y=-x^2-2x+3$
  \task $y=2x^2+4x+3$
  \task $y=-\frac12x^2+2$
\end{tasks}\recorda{$V_x=\dfrac{-b}{2a}$, $V_y=f(V_x)$. Si $a\gt0$, còncava $\cup$; si $a\lt0$, convexa $\cap$.}

\ex{2} Troba la paràbola $y=ax^2+bx+c$ que talla l'eix $OY$ en $(0,3)$, té el vèrtex en $x=2$ i passa pel punt $(1,0)$.

\ex{2} Es llança una pilota cap amunt i la seua altura (en metres) als $t$ segons és $h(t)=-5t^2+20t+1$. Quina altura màxima arriba a tindre i quan? Quan cau a terra?

\ex{2} Troba els punts de tall amb els eixos de $f(x)=x^3-4x^2+3x$ i explica com acaba la gràfica per l'esquerra i per la dreta. Fes-ne un esbós.


% =====================================================
\seccio{3}{Funcions racionals}

\exf{1} Per a cada funció, escriu-la (si cal) com a translació de $y=\frac kx$, indica'n les asímptotes, el domini, el recorregut i la monotonia, i dibuixa-la:\begin{tasks}(4)
  \task $y=\dfrac{2}{x-3}+1$
  \task $y=-\dfrac{3}{x+2}$
  \task $y=\dfrac{2x+1}{x-1}$
  \task $y=\dfrac{x-4}{x-2}$
\end{tasks}

\ex{2} Un grup lloga un autobús per 600~€ i se'n reparteix el preu a parts iguals. Escriu la funció $p(n)$ que dona el que paga cada persona si hi van $n$ persones. Quant paga cadascú si hi van 20 persones? Quantes persones hi van si cadascuna paga 24~€? Què passa amb $p(n)$ quan $n$ es fa molt gran?

\ex{2} Una hipèrbola de la forma $y=\dfrac{k}{x-a}+b$ té les asímptotes $x=2$ i $y=1$ i passa pel punt $(3,4)$. Troba'n l'expressió i el punt de tall amb l'eix $OY$.


% =====================================================
\seccio{4}{Funcions radicals}

\exf{1} Indica el punt on comença cada funció, el domini, el recorregut i la monotonia, i dibuixa-les:\begin{tasks}(4)
  \task $y=\sqrt{x-3}+1$
  \task $y=-\sqrt{x+2}$
  \task $y=\sqrt{4-x}$
  \task $y=-\sqrt{-x-1}+3$
\end{tasks}

\ex{2} Escriu l'expressió de la funció radical representada, que és una translació de $y=\sqrt x$.\begin{center}\begin{tikzpicture}[x=0.364cm,y=0.496cm,baseline={(current bounding box.center)},font=\scriptsize]
\draw[muted!25,line width=.3pt] (-3,-3)--(-3,2);
\draw[muted!25,line width=.3pt] (-2,-3)--(-2,2);
\draw[muted!25,line width=.3pt] (-1,-3)--(-1,2);
\draw[muted!25,line width=.3pt] (0,-3)--(0,2);
\draw[muted!25,line width=.3pt] (1,-3)--(1,2);
\draw[muted!25,line width=.3pt] (2,-3)--(2,2);
\draw[muted!25,line width=.3pt] (3,-3)--(3,2);
\draw[muted!25,line width=.3pt] (4,-3)--(4,2);
\draw[muted!25,line width=.3pt] (5,-3)--(5,2);
\draw[muted!25,line width=.3pt] (6,-3)--(6,2);
\draw[muted!25,line width=.3pt] (7,-3)--(7,2);
\draw[muted!25,line width=.3pt] (8,-3)--(8,2);
\draw[muted!25,line width=.3pt] (-3,-3)--(8,-3);
\draw[muted!25,line width=.3pt] (-3,-2)--(8,-2);
\draw[muted!25,line width=.3pt] (-3,-1)--(8,-1);
\draw[muted!25,line width=.3pt] (-3,0)--(8,0);
\draw[muted!25,line width=.3pt] (-3,1)--(8,1);
\draw[muted!25,line width=.3pt] (-3,2)--(8,2);
\draw[-latex,muted] (-3,0)--(8,0);
\draw[-latex,muted] (0,-3)--(0,2);
\node[below=-1pt,text=muted] at (-3,0) {$-3$};
\node[below=-1pt,text=muted] at (-2,0) {$-2$};
\node[below=-1pt,text=muted] at (-1,0) {$-1$};
\node[below=-1pt,text=muted] at (1,0) {$1$};
\node[below=-1pt,text=muted] at (2,0) {$2$};
\node[below=-1pt,text=muted] at (3,0) {$3$};
\node[below=-1pt,text=muted] at (4,0) {$4$};
\node[below=-1pt,text=muted] at (5,0) {$5$};
\node[below=-1pt,text=muted] at (6,0) {$6$};
\node[below=-1pt,text=muted] at (7,0) {$7$};
\node[below=-1pt,text=muted] at (8,0) {$8$};
\node[left=-1pt,text=muted] at (0,-3) {$-3$};
\node[left=-1pt,text=muted] at (0,-2) {$-2$};
\node[left=-1pt,text=muted] at (0,-1) {$-1$};
\node[left=-1pt,text=muted] at (0,1) {$1$};
\node[left=-1pt,text=muted] at (0,2) {$2$};
\begin{scope}\clip (-3,-3) rectangle (8,2);
\draw[accent,line width=1.1pt,line join=round] (-1,-2)--(-0.928,-1.732)--(-0.856,-1.621)--(-0.784,-1.535)--(-0.712,-1.463)--(-0.64,-1.4)--(-0.568,-1.343)--(-0.496,-1.29)--(-0.424,-1.241)--(-0.352,-1.195)--(-0.28,-1.151)--(-0.208,-1.11)--(-0.136,-1.07)--(-0.064,-1.033)--(0.008,-0.996)--(0.08,-0.961)--(0.152,-0.927)--(0.224,-0.894)--(0.296,-0.862)--(0.368,-0.83)--(0.44,-0.8)--(0.512,-0.77)--(0.584,-0.741)--(0.656,-0.713)--(0.728,-0.685)--(0.8,-0.658)--(0.872,-0.632)--(0.944,-0.606)--(1.016,-0.58)--(1.088,-0.555)--(1.16,-0.53)--(1.232,-0.506)--(1.304,-0.482)--(1.376,-0.459)--(1.448,-0.435)--(1.52,-0.413)--(1.592,-0.39)--(1.664,-0.368)--(1.736,-0.346)--(1.808,-0.324)--(1.88,-0.303)--(1.952,-0.282)--(2.024,-0.261)--(2.096,-0.24)--(2.168,-0.22)--(2.24,-0.2)--(2.312,-0.18)--(2.384,-0.16)--(2.456,-0.141)--(2.528,-0.122)--(2.6,-0.103)--(2.672,-0.084)--(2.744,-0.065)--(2.816,-0.047)--(2.888,-0.028)--(2.96,-0.01)--(3.032,0.008)--(3.104,0.026)--(3.176,0.044)--(3.248,0.061)--(3.32,0.078)--(3.392,0.096)--(3.464,0.113)--(3.536,0.13)--(3.608,0.147)--(3.68,0.163)--(3.752,0.18)--(3.824,0.196)--(3.896,0.213)--(3.968,0.229)--(4.04,0.245)--(4.112,0.261)--(4.184,0.277)--(4.256,0.293)--(4.328,0.308)--(4.4,0.324)--(4.472,0.339)--(4.544,0.355)--(4.616,0.37)--(4.688,0.385)--(4.76,0.4)--(4.832,0.415)--(4.904,0.43)--(4.976,0.445)--(5.048,0.459)--(5.12,0.474)--(5.192,0.488)--(5.264,0.503)--(5.336,0.517)--(5.408,0.531)--(5.48,0.546)--(5.552,0.56)--(5.624,0.574)--(5.696,0.588)--(5.768,0.602)--(5.84,0.615)--(5.912,0.629)--(5.984,0.643)--(6.056,0.656)--(6.128,0.67)--(6.2,0.683)--(6.272,0.697)--(6.344,0.71)--(6.416,0.723)--(6.488,0.736)--(6.56,0.75)--(6.632,0.763)--(6.704,0.776)--(6.776,0.789)--(6.848,0.801)--(6.92,0.814)--(6.992,0.827)--(7.064,0.84)--(7.136,0.852)--(7.208,0.865)--(7.28,0.877)--(7.352,0.89)--(7.424,0.902)--(7.496,0.915)--(7.568,0.927)--(7.64,0.939)--(7.712,0.952)--(7.784,0.964)--(7.856,0.976)--(7.928,0.988)--(8,1);
\filldraw[fill=blue!70!black,draw=blue!70!black,line width=.8pt] (-1,-2) circle (1.5pt);
\filldraw[fill=blue!70!black,draw=blue!70!black,line width=.8pt] (3,0) circle (1.5pt);
\end{scope}\end{tikzpicture}\end{center}

\ex{2} Donada $f(x)=\sqrt{2x-4}+1$, indica on comença, el domini i el recorregut. Per a quin valor de $x$ val $f(x)=3$?


% =====================================================
\seccio{5}{Funcions exponencials i logarítmiques}

\exf{1} Indica l'asímptota horitzontal, el punt de referència (la imatge del $(0,1)$ de $y=a^x$), el recorregut i la monotonia, i dibuixa-les:\begin{tasks}(4)
  \task $y=2^{x-1}+3$
  \task $y=\left(\frac13\right)^x-2$
  \task $y=-3^{x+1}$
  \task $y=1-2^x$
\end{tasks}

\ex{1} Indica l'asímptota vertical, el punt de referència (la imatge del $(1,0)$ de $y=\log_ax$), el domini i la monotonia, i dibuixa-les:\begin{tasks}(4)
  \task $y=\log_2(x+3)$
  \task $y=\log_3x-1$
  \task $y=\log_{\frac12}(x-1)+2$
  \task $y=\log(x-2)+1$
\end{tasks}

\ex{2} La població d'un poble, que ara té 5\,000 habitants, creix un 3\,\% cada any: $P(t)=5\,000\cdot1{,}03^t$. Quants habitants tindrà d'ací a 10 anys? Quant de temps tardarà a duplicar-se?

\ex{2} Un pacient pren 200 mg d'un medicament i cada hora se n'elimina el 20\,\%. Escriu la funció que dona la quantitat que queda al cap de $t$ hores. Quant en queda al cap de 3 hores? Quan quedaran menys de 50 mg?

\ex{2} Les gràfiques són translacions de $y=2^x$ i de $y=\log_2x$. Escriu-ne les expressions.\begin{tasks}(2)
  \task \begin{tikzpicture}[x=0.367cm,y=0.367cm,baseline={(current bounding box.center)},font=\scriptsize]
\draw[muted!25,line width=.3pt] (-5,-4)--(-5,5);
\draw[muted!25,line width=.3pt] (-4,-4)--(-4,5);
\draw[muted!25,line width=.3pt] (-3,-4)--(-3,5);
\draw[muted!25,line width=.3pt] (-2,-4)--(-2,5);
\draw[muted!25,line width=.3pt] (-1,-4)--(-1,5);
\draw[muted!25,line width=.3pt] (0,-4)--(0,5);
\draw[muted!25,line width=.3pt] (1,-4)--(1,5);
\draw[muted!25,line width=.3pt] (2,-4)--(2,5);
\draw[muted!25,line width=.3pt] (3,-4)--(3,5);
\draw[muted!25,line width=.3pt] (4,-4)--(4,5);
\draw[muted!25,line width=.3pt] (-5,-4)--(4,-4);
\draw[muted!25,line width=.3pt] (-5,-3)--(4,-3);
\draw[muted!25,line width=.3pt] (-5,-2)--(4,-2);
\draw[muted!25,line width=.3pt] (-5,-1)--(4,-1);
\draw[muted!25,line width=.3pt] (-5,0)--(4,0);
\draw[muted!25,line width=.3pt] (-5,1)--(4,1);
\draw[muted!25,line width=.3pt] (-5,2)--(4,2);
\draw[muted!25,line width=.3pt] (-5,3)--(4,3);
\draw[muted!25,line width=.3pt] (-5,4)--(4,4);
\draw[muted!25,line width=.3pt] (-5,5)--(4,5);
\draw[-latex,muted] (-5,0)--(4,0);
\draw[-latex,muted] (0,-4)--(0,5);
\node[below=-1pt,text=muted] at (-4,0) {$-4$};
\node[below=-1pt,text=muted] at (-2,0) {$-2$};
\node[below=-1pt,text=muted] at (2,0) {$2$};
\node[below=-1pt,text=muted] at (4,0) {$4$};
\node[left=-1pt,text=muted] at (0,-4) {$-4$};
\node[left=-1pt,text=muted] at (0,-2) {$-2$};
\node[left=-1pt,text=muted] at (0,2) {$2$};
\node[left=-1pt,text=muted] at (0,4) {$4$};
\begin{scope}\clip (-5,-4) rectangle (4,5);
\draw[accent,line width=1.1pt,line join=round] (-5,-2.938)--(-4.928,-2.934)--(-4.856,-2.931)--(-4.784,-2.927)--(-4.712,-2.924)--(-4.64,-2.92)--(-4.568,-2.916)--(-4.496,-2.911)--(-4.424,-2.907)--(-4.352,-2.902)--(-4.28,-2.897)--(-4.208,-2.892)--(-4.136,-2.886)--(-4.064,-2.88)--(-3.992,-2.874)--(-3.92,-2.868)--(-3.848,-2.861)--(-3.776,-2.854)--(-3.704,-2.847)--(-3.632,-2.839)--(-3.56,-2.83)--(-3.488,-2.822)--(-3.416,-2.813)--(-3.344,-2.803)--(-3.272,-2.793)--(-3.2,-2.782)--(-3.128,-2.771)--(-3.056,-2.76)--(-2.984,-2.747)--(-2.912,-2.734)--(-2.84,-2.721)--(-2.768,-2.706)--(-2.696,-2.691)--(-2.624,-2.676)--(-2.552,-2.659)--(-2.48,-2.642)--(-2.408,-2.623)--(-2.336,-2.604)--(-2.264,-2.584)--(-2.192,-2.562)--(-2.12,-2.54)--(-2.048,-2.516)--(-1.976,-2.492)--(-1.904,-2.466)--(-1.832,-2.438)--(-1.76,-2.41)--(-1.688,-2.379)--(-1.616,-2.348)--(-1.544,-2.314)--(-1.472,-2.279)--(-1.4,-2.242)--(-1.328,-2.203)--(-1.256,-2.163)--(-1.184,-2.12)--(-1.112,-2.075)--(-1.04,-2.027)--(-0.968,-1.978)--(-0.896,-1.925)--(-0.824,-1.87)--(-0.752,-1.812)--(-0.68,-1.752)--(-0.608,-1.688)--(-0.536,-1.621)--(-0.464,-1.55)--(-0.392,-1.476)--(-0.32,-1.398)--(-0.248,-1.316)--(-0.176,-1.23)--(-0.104,-1.139)--(-0.032,-1.044)--(0.04,-0.944)--(0.112,-0.839)--(0.184,-0.728)--(0.256,-0.612)--(0.328,-0.489)--(0.4,-0.361)--(0.472,-0.226)--(0.544,-0.084)--(0.616,0.065)--(0.688,0.222)--(0.76,0.387)--(0.832,0.56)--(0.904,0.742)--(0.976,0.934)--(1.048,1.135)--(1.12,1.347)--(1.192,1.569)--(1.264,1.803)--(1.336,2.049)--(1.408,2.307)--(1.48,2.579)--(1.552,2.864)--(1.624,3.165)--(1.696,3.48)--(1.768,3.812)--(1.84,4.16)--(1.912,4.527)--(1.984,4.912)--(2.056,5.317)--(2.128,5.742)--(2.2,6.19)--(2.272,6.66)--(2.344,7.154)--(2.416,7.674)--(2.488,8.22)--(2.56,8.794)--(2.632,9.398)--(2.704,10.032)--(2.776,10.699)--(2.848,11.4)--(2.92,12.137)--(2.992,12.912)--(3.064,13.726)--(3.136,14)--(3.208,14)--(3.28,14)--(3.352,14)--(3.424,14)--(3.496,14)--(3.568,14)--(3.64,14)--(3.712,14)--(3.784,14)--(3.856,14)--(3.928,14)--(4,14);
\draw[muted,dashed,line width=.5pt] (-5,-3)--(4,-3);
\filldraw[fill=blue!70!black,draw=blue!70!black,line width=.8pt] (-1,-2) circle (1.5pt);
\filldraw[fill=blue!70!black,draw=blue!70!black,line width=.8pt] (1,1) circle (1.5pt);
\end{scope}\end{tikzpicture}
  \task \begin{tikzpicture}[x=0.367cm,y=0.367cm,baseline={(current bounding box.center)},font=\scriptsize]
\draw[muted!25,line width=.3pt] (-1,-3)--(-1,4);
\draw[muted!25,line width=.3pt] (0,-3)--(0,4);
\draw[muted!25,line width=.3pt] (1,-3)--(1,4);
\draw[muted!25,line width=.3pt] (2,-3)--(2,4);
\draw[muted!25,line width=.3pt] (3,-3)--(3,4);
\draw[muted!25,line width=.3pt] (4,-3)--(4,4);
\draw[muted!25,line width=.3pt] (5,-3)--(5,4);
\draw[muted!25,line width=.3pt] (6,-3)--(6,4);
\draw[muted!25,line width=.3pt] (7,-3)--(7,4);
\draw[muted!25,line width=.3pt] (8,-3)--(8,4);
\draw[muted!25,line width=.3pt] (-1,-3)--(8,-3);
\draw[muted!25,line width=.3pt] (-1,-2)--(8,-2);
\draw[muted!25,line width=.3pt] (-1,-1)--(8,-1);
\draw[muted!25,line width=.3pt] (-1,0)--(8,0);
\draw[muted!25,line width=.3pt] (-1,1)--(8,1);
\draw[muted!25,line width=.3pt] (-1,2)--(8,2);
\draw[muted!25,line width=.3pt] (-1,3)--(8,3);
\draw[muted!25,line width=.3pt] (-1,4)--(8,4);
\draw[-latex,muted] (-1,0)--(8,0);
\draw[-latex,muted] (0,-3)--(0,4);
\node[below=-1pt,text=muted] at (2,0) {$2$};
\node[below=-1pt,text=muted] at (4,0) {$4$};
\node[below=-1pt,text=muted] at (6,0) {$6$};
\node[below=-1pt,text=muted] at (8,0) {$8$};
\node[left=-1pt,text=muted] at (0,-2) {$-2$};
\node[left=-1pt,text=muted] at (0,2) {$2$};
\node[left=-1pt,text=muted] at (0,4) {$4$};
\begin{scope}\clip (-1,-3) rectangle (8,4);
\draw[accent,line width=1.1pt,line join=round] (1.005,-6.644)--(1.061,-3.036)--(1.117,-2.096)--(1.173,-1.532)--(1.229,-1.128)--(1.285,-0.812)--(1.341,-0.553)--(1.397,-0.334)--(1.453,-0.143)--(1.509,0.025)--(1.565,0.175)--(1.621,0.312)--(1.677,0.436)--(1.732,0.551)--(1.788,0.657)--(1.844,0.756)--(1.9,0.849)--(1.956,0.936)--(2.012,1.018)--(2.068,1.095)--(2.124,1.169)--(2.18,1.239)--(2.236,1.306)--(2.292,1.37)--(2.348,1.431)--(2.404,1.49)--(2.46,1.546)--(2.516,1.6)--(2.572,1.652)--(2.628,1.703)--(2.684,1.752)--(2.74,1.799)--(2.796,1.845)--(2.852,1.889)--(2.908,1.932)--(2.964,1.974)--(3.02,2.014)--(3.076,2.053)--(3.131,2.092)--(3.187,2.129)--(3.243,2.166)--(3.299,2.201)--(3.355,2.236)--(3.411,2.27)--(3.467,2.303)--(3.523,2.335)--(3.579,2.367)--(3.635,2.398)--(3.691,2.428)--(3.747,2.458)--(3.803,2.487)--(3.859,2.515)--(3.915,2.543)--(3.971,2.571)--(4.027,2.598)--(4.083,2.624)--(4.139,2.65)--(4.195,2.676)--(4.251,2.701)--(4.307,2.725)--(4.363,2.75)--(4.419,2.773)--(4.475,2.797)--(4.53,2.82)--(4.586,2.843)--(4.642,2.865)--(4.698,2.887)--(4.754,2.909)--(4.81,2.93)--(4.866,2.951)--(4.922,2.972)--(4.978,2.992)--(5.034,3.012)--(5.09,3.032)--(5.146,3.052)--(5.202,3.071)--(5.258,3.09)--(5.314,3.109)--(5.37,3.128)--(5.426,3.146)--(5.482,3.164)--(5.538,3.182)--(5.594,3.2)--(5.65,3.217)--(5.706,3.234)--(5.762,3.251)--(5.818,3.268)--(5.874,3.285)--(5.929,3.301)--(5.985,3.318)--(6.041,3.334)--(6.097,3.35)--(6.153,3.366)--(6.209,3.381)--(6.265,3.396)--(6.321,3.412)--(6.377,3.427)--(6.433,3.442)--(6.489,3.457)--(6.545,3.471)--(6.601,3.486)--(6.657,3.5)--(6.713,3.514)--(6.769,3.528)--(6.825,3.542)--(6.881,3.556)--(6.937,3.57)--(6.993,3.583)--(7.049,3.597)--(7.105,3.61)--(7.161,3.623)--(7.217,3.636)--(7.273,3.649)--(7.328,3.662)--(7.384,3.675)--(7.44,3.687)--(7.496,3.7)--(7.552,3.712)--(7.608,3.724)--(7.664,3.736)--(7.72,3.749)--(7.776,3.76)--(7.832,3.772)--(7.888,3.784)--(7.944,3.796)--(8,3.807);
\draw[muted,dashed,line width=.5pt] (1,-3)--(1,4);
\filldraw[fill=blue!70!black,draw=blue!70!black,line width=.8pt] (2,1) circle (1.5pt);
\filldraw[fill=blue!70!black,draw=blue!70!black,line width=.8pt] (3,2) circle (1.5pt);
\end{scope}\end{tikzpicture}
\end{tasks}


% =====================================================
\seccio{6}{Funcions trigonomètriques}

\exf{1} Respon fent servir les gràfiques de les funcions trigonomètriques:\begin{tasks}(1)
  \task En quins punts de $[0,2\pi]$ talla l'eix $OX$ la funció $y=\cos x$? On té els màxims i els mínims en $[0,4\pi]$?
  \task Quin és el domini de $y=\tan x$? I el període?
  \task Indica el recorregut i el període de $y=3\sin x-1$ i de $y=2+\cos x$.
  \task Estudia la simetria de $f(x)=\sin x\cdot\cos x$ i de $g(x)=x\sin x$.
\end{tasks}

\ex{2} Dibuixa en $[0,2\pi]$ les funcions $y=\sin x+1$ i $y=2\cos x$ i indica'n el recorregut, els màxims, els mínims i els talls amb l'eix $OX$.


% =====================================================
\seccio{7}{Funcions a trossos i valor absolut}

\exf{1} Representa $f(x)=\begin{cases}2x+4, & \text{si } x\lt-1\\ x^2-1, & \text{si } -1\le x\lt2\\ 3, & \text{si } x\geq2\end{cases}$, calcula $f(-3)$, $f(0)$, $f(2)$ i $f(5)$, i estudia'n la continuïtat i el recorregut.

\ex{2} Escriu l'expressió de la funció a trossos representada:\begin{center}\begin{tikzpicture}[x=0.467cm,y=0.467cm,baseline={(current bounding box.center)},font=\scriptsize]
\draw[muted!25,line width=.3pt] (-4,-1)--(-4,5);
\draw[muted!25,line width=.3pt] (-3,-1)--(-3,5);
\draw[muted!25,line width=.3pt] (-2,-1)--(-2,5);
\draw[muted!25,line width=.3pt] (-1,-1)--(-1,5);
\draw[muted!25,line width=.3pt] (0,-1)--(0,5);
\draw[muted!25,line width=.3pt] (1,-1)--(1,5);
\draw[muted!25,line width=.3pt] (2,-1)--(2,5);
\draw[muted!25,line width=.3pt] (3,-1)--(3,5);
\draw[muted!25,line width=.3pt] (4,-1)--(4,5);
\draw[muted!25,line width=.3pt] (5,-1)--(5,5);
\draw[muted!25,line width=.3pt] (-4,-1)--(5,-1);
\draw[muted!25,line width=.3pt] (-4,0)--(5,0);
\draw[muted!25,line width=.3pt] (-4,1)--(5,1);
\draw[muted!25,line width=.3pt] (-4,2)--(5,2);
\draw[muted!25,line width=.3pt] (-4,3)--(5,3);
\draw[muted!25,line width=.3pt] (-4,4)--(5,4);
\draw[muted!25,line width=.3pt] (-4,5)--(5,5);
\draw[-latex,muted] (-4,0)--(5,0);
\draw[-latex,muted] (0,-1)--(0,5);
\node[below=-1pt,text=muted] at (-4,0) {$-4$};
\node[below=-1pt,text=muted] at (-3,0) {$-3$};
\node[below=-1pt,text=muted] at (-2,0) {$-2$};
\node[below=-1pt,text=muted] at (-1,0) {$-1$};
\node[below=-1pt,text=muted] at (1,0) {$1$};
\node[below=-1pt,text=muted] at (2,0) {$2$};
\node[below=-1pt,text=muted] at (3,0) {$3$};
\node[below=-1pt,text=muted] at (4,0) {$4$};
\node[below=-1pt,text=muted] at (5,0) {$5$};
\node[left=-1pt,text=muted] at (0,-1) {$-1$};
\node[left=-1pt,text=muted] at (0,1) {$1$};
\node[left=-1pt,text=muted] at (0,2) {$2$};
\node[left=-1pt,text=muted] at (0,3) {$3$};
\node[left=-1pt,text=muted] at (0,4) {$4$};
\node[left=-1pt,text=muted] at (0,5) {$5$};
\begin{scope}\clip (-4,-1) rectangle (5,5);
\draw[accent,line width=1.1pt,line join=round] (-4,4)--(-3.968,3.968)--(-3.936,3.936)--(-3.904,3.904)--(-3.872,3.872)--(-3.84,3.84)--(-3.808,3.808)--(-3.776,3.776)--(-3.744,3.744)--(-3.712,3.712)--(-3.68,3.68)--(-3.648,3.648)--(-3.616,3.616)--(-3.584,3.584)--(-3.552,3.552)--(-3.52,3.52)--(-3.488,3.488)--(-3.456,3.456)--(-3.424,3.424)--(-3.392,3.392)--(-3.36,3.36)--(-3.328,3.328)--(-3.296,3.296)--(-3.264,3.264)--(-3.232,3.232)--(-3.2,3.2)--(-3.168,3.168)--(-3.136,3.136)--(-3.104,3.104)--(-3.072,3.072)--(-3.04,3.04)--(-3.008,3.008)--(-2.976,2.976)--(-2.944,2.944)--(-2.912,2.912)--(-2.88,2.88)--(-2.848,2.848)--(-2.816,2.816)--(-2.784,2.784)--(-2.752,2.752)--(-2.72,2.72)--(-2.688,2.688)--(-2.656,2.656)--(-2.624,2.624)--(-2.592,2.592)--(-2.56,2.56)--(-2.528,2.528)--(-2.496,2.496)--(-2.464,2.464)--(-2.432,2.432)--(-2.4,2.4)--(-2.368,2.368)--(-2.336,2.336)--(-2.304,2.304)--(-2.272,2.272)--(-2.24,2.24)--(-2.208,2.208)--(-2.176,2.176)--(-2.144,2.144)--(-2.112,2.112)--(-2.08,2.08)--(-2.048,2.048)--(-2.016,2.016)--(-1.984,1.984)--(-1.952,1.952)--(-1.92,1.92)--(-1.888,1.888)--(-1.856,1.856)--(-1.824,1.824)--(-1.792,1.792)--(-1.76,1.76)--(-1.728,1.728)--(-1.696,1.696)--(-1.664,1.664)--(-1.632,1.632)--(-1.6,1.6)--(-1.568,1.568)--(-1.536,1.536)--(-1.504,1.504)--(-1.472,1.472)--(-1.44,1.44)--(-1.408,1.408)--(-1.376,1.376)--(-1.344,1.344)--(-1.312,1.312)--(-1.28,1.28)--(-1.248,1.248)--(-1.216,1.216)--(-1.184,1.184)--(-1.152,1.152)--(-1.12,1.12)--(-1.088,1.088)--(-1.056,1.056)--(-1.024,1.024)--(-0.992,0.992)--(-0.96,0.96)--(-0.928,0.928)--(-0.896,0.896)--(-0.864,0.864)--(-0.832,0.832)--(-0.8,0.8)--(-0.768,0.768)--(-0.736,0.736)--(-0.704,0.704)--(-0.672,0.672)--(-0.64,0.64)--(-0.608,0.608)--(-0.576,0.576)--(-0.544,0.544)--(-0.512,0.512)--(-0.48,0.48)--(-0.448,0.448)--(-0.416,0.416)--(-0.384,0.384)--(-0.352,0.352)--(-0.32,0.32)--(-0.288,0.288)--(-0.256,0.256)--(-0.224,0.224)--(-0.192,0.192)--(-0.16,0.16)--(-0.128,0.128)--(-0.096,0.096)--(-0.064,0.064)--(-0.032,0.032)--(0,-0);
\draw[accent,line width=1.1pt,line join=round] (0,0)--(0.016,0)--(0.032,0.001)--(0.048,0.002)--(0.064,0.004)--(0.08,0.006)--(0.096,0.009)--(0.112,0.013)--(0.128,0.016)--(0.144,0.021)--(0.16,0.026)--(0.176,0.031)--(0.192,0.037)--(0.208,0.043)--(0.224,0.05)--(0.24,0.058)--(0.256,0.066)--(0.272,0.074)--(0.288,0.083)--(0.304,0.092)--(0.32,0.102)--(0.336,0.113)--(0.352,0.124)--(0.368,0.135)--(0.384,0.147)--(0.4,0.16)--(0.416,0.173)--(0.432,0.187)--(0.448,0.201)--(0.464,0.215)--(0.48,0.23)--(0.496,0.246)--(0.512,0.262)--(0.528,0.279)--(0.544,0.296)--(0.56,0.314)--(0.576,0.332)--(0.592,0.35)--(0.608,0.37)--(0.624,0.389)--(0.64,0.41)--(0.656,0.43)--(0.672,0.452)--(0.688,0.473)--(0.704,0.496)--(0.72,0.518)--(0.736,0.542)--(0.752,0.566)--(0.768,0.59)--(0.784,0.615)--(0.8,0.64)--(0.816,0.666)--(0.832,0.692)--(0.848,0.719)--(0.864,0.746)--(0.88,0.774)--(0.896,0.803)--(0.912,0.832)--(0.928,0.861)--(0.944,0.891)--(0.96,0.922)--(0.976,0.953)--(0.992,0.984)--(1.008,1.016)--(1.024,1.049)--(1.04,1.082)--(1.056,1.115)--(1.072,1.149)--(1.088,1.184)--(1.104,1.219)--(1.12,1.254)--(1.136,1.29)--(1.152,1.327)--(1.168,1.364)--(1.184,1.402)--(1.2,1.44)--(1.216,1.479)--(1.232,1.518)--(1.248,1.558)--(1.264,1.598)--(1.28,1.638)--(1.296,1.68)--(1.312,1.721)--(1.328,1.764)--(1.344,1.806)--(1.36,1.85)--(1.376,1.893)--(1.392,1.938)--(1.408,1.982)--(1.424,2.028)--(1.44,2.074)--(1.456,2.12)--(1.472,2.167)--(1.488,2.214)--(1.504,2.262)--(1.52,2.31)--(1.536,2.359)--(1.552,2.409)--(1.568,2.459)--(1.584,2.509)--(1.6,2.56)--(1.616,2.611)--(1.632,2.663)--(1.648,2.716)--(1.664,2.769)--(1.68,2.822)--(1.696,2.876)--(1.712,2.931)--(1.728,2.986)--(1.744,3.042)--(1.76,3.098)--(1.776,3.154)--(1.792,3.211)--(1.808,3.269)--(1.824,3.327)--(1.84,3.386)--(1.856,3.445)--(1.872,3.504)--(1.888,3.565)--(1.904,3.625)--(1.92,3.686)--(1.936,3.748)--(1.952,3.81)--(1.968,3.873)--(1.984,3.936)--(2,4);
\draw[accent,line width=1.1pt,line join=round] (2,1)--(5,1);
\filldraw[fill=accent,draw=accent,line width=.8pt] (2,4) circle (1.5pt);
\filldraw[fill=white,draw=accent,line width=.8pt] (2,1) circle (1.5pt);
\end{scope}\end{tikzpicture}\end{center}

\ex{1} Escriu a trossos i representa:\begin{tasks}(4)
  \task $y=|x-3|$
  \task $y=|x^2-4x|$
  \task $y=|x^2-x-2|$
  \task $y=|2x+1|$
\end{tasks}\recorda{troba on s'anul·la l'expressió de dins, estudia'n el signe i canvia-li el signe on és negativa.}

\ex{2} Un aparcament cobra 2~€ per la primera hora (o fracció) i $1{,}50$~€ per cada hora de més. Escriu la funció a trossos que dona el preu segons les hores $x$ (amb $x\gt0$, comptant les hores de manera contínua), calcula el preu de 3 hores i quantes hores has estat si pagues 8~€. És contínua en $x=1$?

\ex{2} Resol gràficament l'equació $|x-1|=3$: dibuixa $y=|x-1|$ i la recta $y=3$, i troba on es tallen.


% =====================================================
\seccio{8}{Problemes amb funcions quadràtiques}

\exf{1} Troba el màxim i el mínim de $f(x)=x^2-4x+1$ en l'interval $[-1,3]$ i en l'interval $[3,5]$.\recorda{en un interval tancat, el màxim i el mínim d'una paràbola són al vèrtex (si està dins de l'interval) o en els extrems de l'interval. Calcula'ls tots i compara.}

\ex{2} Dos nombres sumen 20. Quins són si el seu producte és el més gran possible? I si la suma dels seus quadrats és la més petita possible?

\ex{2} Una botiga ven un producte a $p$ euros i, a aquest preu, en ven $200-4p$ unitats al mes ($0\le p\le50$). Escriu la funció d'ingressos $I(p)$. A quin preu són màxims els ingressos i quant valen? Si per motius de competència el preu ha d'estar entre 30 i 45~€, quin preu convé posar?

\ex{2} Una empresa fabrica $x$ unitats d'un article ($0\le x\le50$). Els ingressos són $I(x)=60x$ i els costos $C(x)=x^2+20x+300$ (en euros). Escriu la funció de beneficis $B(x)=I(x)-C(x)$. Quantes unitats ha de fabricar per a obtindre el benefici màxim? Entre quantes unitats hi ha guanys? Si només pot fabricar-ne 15 com a màxim, quin és el benefici més gran?

\ex{2} El cost per unitat (en euros) de fabricar $x$ centenars de peces és $c(x)=0{,}5x^2-20x+250$, amb $0\le x\le40$. Quina producció fa mínim el cost per unitat? Quin és aquest cost? I si la màquina només pot fer com a molt 1\,500 peces?

\ex{2} Amb 60 m de tanca volem fer un corral rectangular arrimat a una paret (el costat de la paret no es tanca). Si els dos costats perpendiculars a la paret mesuren $x$, escriu l'àrea $A(x)$ i troba les dimensions del corral d'àrea màxima.

\ex{3} Un hotel té 40 habitacions. Si cobra 80~€ per nit, les ompli totes; per cada 5~€ que puja el preu, queda una habitació buida. Si el puja $n$ vegades 5~€, escriu els ingressos $I(n)$. Quin preu fa màxims els ingressos? Quantes habitacions s'ocupen?

\vfill
\begin{center}\color{muted}\small\faCheckCircle[regular]\ \ Ànim! Comprova els teus resultats amb les eines interactives dels apunts de la unitat.\end{center}

\end{document}
